\section{Nicht-funktionale Anforderungen}
\label{sec:nfr}

\subsection{Architektur}
\label{sec:nfr_architektur}

\textbf{NFR-ARC-01}\\
Die Arduinos müssen ROS-kompatible Software ausführen (z. B. micro-ROS), sodass sie vollständig in die bestehende ROS1/ROS2-basierten Systemarchitektur integriert sind.

\textbf{NFR-ARC-02}\\
Alle Nothalt-relevanten Meldungen müssen als ROS-Nachrichten veröffentlicht werden (z. B. \texttt{/emergency\_stop\_state}, \texttt{/fault\_state}, \texttt{/sensor\_health}).


% =====================================================================
\subsection{Zeitverhalten \& Performance}
\label{sec:nfr_time}

\textbf{NFR-TIM-01}\\
Zwischen der Erkennung eines Hindernisses in der kritischen Zone und dem Versand des Nothaltkommandos an den Master-Arduino dürfen maximal 50\,ms vergehen.

\textbf{NFR-TIM-02}\\
Das System muss sicherstellen, dass das Fahrzeug bei einer Geschwindigkeit von bis zu 0{,}55\,m/s vor einem Hindernis zum Stillstand kommt und dabei einen Sicherheitsabstand von mindestens 20\,cm einhält.

\textbf{NFR-TIM-03}\\
Die Sensorverarbeitung muss mindestens folgende Raten unterstützen:
\begin{itemize}
    \item 2D-Laserscanner: 40\,Hz,
    \item 3D-Lidar: 20--30\,Hz,
    \item Encoder: 50\,Hz.
\end{itemize}


% =====================================================================
\subsection{Sicherheit}
\label{sec:nfr_safety}

\textbf{NFR-SAF-01}\\
Im Falle widersprüchlicher oder fehlender Sensordaten muss das System sicher reagieren, indem es entweder in den Warnmodus wechselt oder einen Nothalt durchführt (je nach Situation).

\textbf{NFR-SAF-02}\\
Die \ac{Hardstop}-Funktion muss unabhängig von der Software zuverlässig funktionieren.


% =====================================================================
\subsection{Zuverlässigkeit}
\label{sec:nfr_reliability}

\textbf{NFR-REL-01}\\
Das System soll so ausgelegt sein, dass Fehlbremsungen (False Positives) höchstens 5\,\% betragen.

\textbf{NFR-REL-02}\\
Nicht erkannte relevante Hindernisse (False Negatives) dürfen höchstens 1\,\% betragen.


% =====================================================================
\subsection{Logging}
\label{sec:nfr_logging}

\textbf{NFR-LOG-01}\\
Das Logging-System muss ROS-Bag-Dateien als primäres Speicherformat verwenden.

\textbf{NFR-LOG-02}\\
Logdaten müssen so lange gespeichert werden, bis
\begin{itemize}
    \item der Benutzer sie manuell löscht oder
    \item sie durch eine automatische Aufräumroutine gelöscht werden.
\end{itemize}
Die automatische Aufräumroutine muss Logdateien löschen, die älter als 14 Tage sind.

\textbf{NFR-LOG-03}\\
Arduinos müssen Sensordaten und Zustandsdaten über ROS-Topics bereitstellen, sodass sie automatisch in ROS-Bags gespeichert werden können.
